% BEGIN PROFESSOR FEEDBACK PF14: Implementation - Introduction and libraries
\section{Implementation details}
\label{sec:implementation}

This section records the tools the pipeline is built on, how the code is
organised, and the measures that make a run reproducible. The engineering
challenges are documented in Appendix~\ref{app:engineering_challenges}.

\subsection{Language and libraries}

The pipeline is written in Python, with the temporal motif counting carried out by
Raphtory's Rust extension. Everything runs on the CPU, so no graphics card is
needed, and Raphtory, DuckDB, the random forest, and XGBoost are all set up to use
the thread allocation recorded for the run. Native thread limits are set before
the numerical libraries are imported, so two coordinated workers cannot each
claim every core. Development and validation used
Python~3.13 in an isolated virtual environment. The exact interpreter and package
versions of each run are recorded in its protocol, and that record, rather than the
environment described here, is what counts for the final numbers.
Table~\ref{tab:libraries} lists the main libraries and their role.

\begin{table}[h]
\centering
\caption{The main libraries used by the pipeline.}
\label{tab:libraries}
\renewcommand{\arraystretch}{1.3}
\small
\begin{tabular}{p{2.7cm} p{1.5cm} p{8.5cm}}
\toprule
\textbf{Library} & \textbf{Version} & \textbf{Role in the pipeline} \\
\midrule

\textbf{Raphtory} &
  0.17.0 &
  Stores the temporal multigraph and runs the parallel Rust routine
  \texttt{local\_temporal\_three\_node\_motifs}, using the local counting
  approach of \citet{arnold2024motifs}. Its full 40-role output is retained for HuggingFace.
  The 32 star and two-node roles are retained for each Real-CATS layer. \\

\textbf{DuckDB} &
  1.5.4 &
  Queries the prepared Parquet directly, computes the baseline, and sorts events
  for each ordering variant with bounded memory. \\

\textbf{PyArrow} &
  23.0.1 &
  Streams the full HuggingFace edge source, writes the prepared Parquet files,
  and carries the ordered batches from DuckDB to Raphtory. \\

\textbf{pandas} &
  3.0.2 &
  Holds the address-level targets, fingerprints, model matrices, predictions, and
  metrics. The primary HF and Real-CATS event tables are not materialised
  together in pandas. The much smaller BLTE table is. \\

\textbf{NumPy} &
  2.4.4 &
  Provides numerical arrays and the deterministic random numbers used by the
  models, the resampling, and the metrics. \\

\textbf{scikit-learn} &
  1.8.0 &
  Supplies the group-aware and stratified splitting, the random forest, the
  standardisation and logistic regression, and the metrics. \\

\textbf{XGBoost} &
  3.2.0 &
  Provides the gradient-boosted robustness model for the combined RQ1 view. \\

\textbf{matplotlib} &
  3.10.8 &
  Produces the figures written for the Results chapter. \\

\bottomrule
\end{tabular}
\end{table}

NetworkX and the HuggingFace \texttt{datasets} client are used only by the
exploratory application and by the data download, not by the locked pipeline.
SQLite, which comes with Python, provides the index used to de-duplicate the
Real-CATS histories.
% END PROFESSOR FEEDBACK PF14
